\section{Αποδοτική Προσαρμογή Παραμέτρων (\en{Parameter-Efficient Fine-Tuning}) και \en{Low-Rank Adaptation}}

\subsection{Εισαγωγή}

Η ραγδαία αύξηση του μεγέθους των Μεγάλων Γλωσσικών Μοντέλων (\en{LLMs}) έχει οδηγήσει σε σημαντικές βελτιώσεις απόδοσης σε πλήθος εφαρμογών. Ωστόσο, αυτές οι βελτιώσεις συνοδεύονται από αυξημένες απαιτήσεις σε υπολογιστικούς πόρους και μνήμη. Τα σύγχρονα μοντέλα περιέχουν δισεκατομμύρια παραμέτρους, καθιστώντας την πλήρη προσαρμογή (\en{fine-tuning}) μη πρακτική σε πολλά σενάρια, ιδιαίτερα σε περιβάλλοντα με περιορισμένους πόρους, όπως το \en{edge computing} ή κατανεμημένα δίκτυα \en{5G}/\en{6G}.

Για την αντιμετώπιση αυτών των προκλήσεων, έχουν προταθεί τεχνικές \en{Parameter-Efficient Fine-Tuning (PEFT)}. Οι μέθοδοι αυτές επιτρέπουν την προσαρμογή προεκπαιδευμένων μοντέλων ενημερώνοντας μόνο ένα μικρό υποσύνολο παραμέτρων ή εισάγοντας ελαφριά εκπαιδεύσιμα στοιχεία. Μεταξύ αυτών, η τεχνική \en{Low-Rank Adaptation (LoRA)} έχει αναδειχθεί ως μία από τις πιο αποδοτικές και ευρέως χρησιμοποιούμενες τεχνικές \cite{hu2021lora}.

\subsection{Κίνητρο για \en{Parameter-Efficient Fine-Tuning}}

Η παραδοσιακή προσαρμογή περιλαμβάνει την ενημέρωση όλων των παραμέτρων ενός μοντέλου. Παρότι επιτυγχάνει υψηλή ακρίβεια, παρουσιάζει σημαντικούς περιορισμούς.

Πρώτον, το υπολογιστικό κόστος είναι ιδιαίτερα υψηλό και απαιτεί εξειδικευμένο υλικό όπως \en{GPUs} ή \en{TPUs}. Δεύτερον, η αποθήκευση πολλαπλών εκδόσεων μοντέλων για διαφορετικές εργασίες είναι αναποτελεσματική. Τρίτον, υπάρχει κίνδυνος υπερπροσαρμογής και φαινόμενα όπως \en{catastrophic forgetting} \cite{raffel2020exploring}.

Οι τεχνικές \en{PEFT} αντιμετωπίζουν αυτά τα προβλήματα επιτρέποντας αποδοτική προσαρμογή με ελάχιστες αλλαγές στις παραμέτρους.

\subsection{Επισκόπηση Τεχνικών \en{PEFT}}

\subsubsection{\en{Adapters}}

Τα \en{adapters} είναι μικρές μονάδες που ενσωματώνονται σε προεκπαιδευμένα μοντέλα, συνήθως μεταξύ των επιπέδων ενός \en{Transformer}. Εκπαιδεύονται μόνο αυτές, ενώ οι αρχικές παράμετροι παραμένουν σταθερές \cite{houlsby2019parameter}.

\subsubsection{\en{Prompt Tuning}}

Το \en{prompt tuning} τροποποιεί την είσοδο του μοντέλου μέσω εκπαιδεύσιμων ενσωματώσεων (\en{soft prompts}) \cite{lester2021power}. Είναι ιδιαίτερα αποδοτικό, αλλά ευαίσθητο στον σχεδιασμό των \en{prompts}.

\subsubsection{\en{Low-Rank Adaptation (LoRA)}}

Η \en{LoRA} αποτελεί τεχνική που τροποποιεί απευθείας τα βάρη του μοντέλου με αποδοτικό τρόπο, χωρίς να τα επανεκπαιδεύει πλήρως.

\subsection{\en{Low-Rank Adaptation}: Βασική Ιδέα}

Η βασική ιδέα της \en{LoRA} είναι η προσέγγιση των ενημερώσεων βαρών μέσω αποσύνθεσης χαμηλού βαθμού \cite{hu2021lora}.

\begin{itemize}
    \item \( A \in \mathbb{R}^{d \times r} \)
    \item \( B \in \mathbb{R}^{r \times k} \)
    \item \( r \ll \min(d,k) \)
    \item \( \Delta W = A B \in \mathbb{R}^{d \times k} \)
\end{itemize}

Το τελικό βάρος είναι:

\[
W' = W + \Delta W
\]

Μόνο οι πίνακες \(A\) και \(B\) εκπαιδεύονται, μειώνοντας δραστικά τον αριθμό παραμέτρων \cite{hu2021lora}.

\subsection{Εννοιολογική Εξήγηση Προσέγγισης Χαμηλού Βαθμού}

Η προσέγγιση χαμηλού βαθμού επιτρέπει την αναπαράσταση μεγάλων πινάκων με μικρότερους, διατηρώντας την πιο σημαντική πληροφορία. Στα νευρωνικά δίκτυα, πολλές παράμετροι είναι πλεονασματικές και μπορούν να συμπιεστούν χωρίς σημαντική απώλεια απόδοσης.

Η τεχνική \en{LoRA} εκμεταλλεύεται αυτή την ιδιότητα, καθιστώντας την ιδιαίτερα αποτελεσματική σε μοντέλα \en{Transformer}.

\subsection{Πλεονεκτήματα της \en{LoRA}}

\subsubsection{Μείωση Εκπαιδεύσιμων Παραμέτρων}

Μειώνει δραστικά τις απαιτήσεις εκπαίδευσης, επιτρέποντας χρήση ακόμη και σε περιορισμένο υλικό.

\subsubsection{Μικρότερη Κατανάλωση Μνήμης}

Απαιτεί λιγότερη μνήμη, γεγονός σημαντικό για \en{edge} και κατανεμημένα περιβάλλοντα.

\subsubsection{Ταχύτερη Εκπαίδευση}

Η εκπαίδευση είναι ταχύτερη λόγω μειωμένης πολυπλοκότητας.

\subsubsection{Καταλληλότητα για Κατανεμημένα Συστήματα}

Η \en{LoRA} είναι ιδανική για περιβάλλοντα \en{5G}/\en{6G} λόγω χαμηλού κόστους και επεκτασιμότητας.

\subsection{Συμπέρασμα}

Οι τεχνικές \en{PEFT} προσφέρουν πρακτική λύση στα προβλήματα προσαρμογής μεγάλων μοντέλων. H τεχνική \en{LoRA} ξεχωρίζει ως αποδοτική και ισχυρή προσέγγιση.

Στο πλαίσιο της παρούσας εργασίας, αποτελεί κατάλληλο εργαλείο για την εφαρμογή \en{LLMs} στην ανάλυση αρχείων καταγραφής και την κυβερνοασφάλεια, ιδιαίτερα σε περιβάλλοντα όπου η αποδοτικότητα και η επεκτασιμότητα είναι κρίσιμες.  